\documentclass[11pt,a4paper]{article}

\usepackage[utf8]{inputenc}
\usepackage[T1]{fontenc}
\usepackage[french]{babel}
\usepackage[a4paper,margin=2.1cm]{geometry}
\usepackage{amsmath,amssymb,amsthm}
\usepackage{lmodern}
\usepackage{enumitem}
\usepackage{hyperref}
\usepackage{fancyhdr}
\usepackage{xcolor}
\usepackage{tcolorbox}
\usepackage{titlesec}
\usepackage{mathrsfs}
\usepackage{tikz}
\usepackage{pgfplots}
\usepackage{array}
\usepackage{booktabs}
\usepackage{multicol}
\usepackage{caption}

\pgfplotsset{compat=1.18}
\usetikzlibrary{arrows.meta,calc,patterns,decorations.pathreplacing}

\hypersetup{
    colorlinks=true,
    linkcolor=black,
    urlcolor=blue
}

\setlength{\parindent}{0pt}
\setlength{\parskip}{0.55em}

\pagestyle{fancy}
\fancyhf{}
\fancyhead[L]{DM de Physique-Chimie}
\fancyhead[R]{Terminale générale}
\fancyfoot[C]{\thepage}

\titleformat{\section}{\large\bfseries}{Exercice \thesection{} --}{0.6em}{}

\newtcolorbox{consignesbox}{
    colback=gray!8,
    colframe=black,
    boxrule=0.8pt,
    arc=2mm
}

\newtcolorbox{objectifbox}{
    colback=gray!5,
    colframe=black,
    boxrule=0.6pt,
    arc=1.5mm
}

\newtcolorbox{schemabox}{
    colback=white,
    colframe=black,
    boxrule=0.5pt,
    arc=1.2mm
}

\newcommand{\R}{\mathbb{R}}
\newcommand{\dd}{\,\mathrm{d}}

\begin{document}

\begin{center}
    {\LARGE \textbf{Devoir Maison de Physique-Chimie}}\\[0.3em]
    {\large \textbf{Niveau Terminale générale}}\\[0.3em]
    Durée conseillée : \textbf{4 h}
\end{center}

\begin{consignesbox}
\textbf{Consignes pour le DM :}
\begin{itemize}[leftmargin=1.2cm]
    \item proposer une rédaction de type DS, soignée et claire ;
    \item faire tous les exercices ; si vous bloquez plus de 30 minutes sur un exercice, passez au suivant ;
    \item ne pas utiliser ChatGPT ni aucune aide extérieure ;
    \item noter dans la marge le temps que vous avez consacré au DM.
\end{itemize}
\end{consignesbox}

\begin{objectifbox}
\textbf{Esprit du sujet.}  
Ce devoir vise à tester la compréhension des chapitres essentiels du programme de Terminale, avec une exigence de rédaction et de raisonnement de type DS.
\end{objectifbox}

%===========================================================
\section{Titrage et validité d’un protocole}

On considère une solution aqueuse contenant un acide faible monoprotique noté \(AH\).  
On souhaite déterminer sa concentration par titrage à l’aide d’une solution d’hydroxyde de sodium.

La réaction support du titrage est :
\[
AH_{(aq)} + HO^-_{(aq)} = A^-_{(aq)} + H_2O_{(l)}.
\]

On prélève un volume \(V_A\) de solution acide, puis on le titre par une solution basique de concentration \(C_B\). Le volume versé à l’équivalence est noté \(V_E\).



\begin{enumerate}[label=\textbf{\arabic*.}]
    \item Définir précisément l’équivalence dans un titrage.

    \item Montrer qu’à l’équivalence, la concentration \(C_A\) de la solution acide vérifie :
    \[
    C_A=\frac{C_BV_E}{V_A}.
    \]
    \item Un élève affirme : \og si on double le volume prélevé \(V_A\), alors le volume équivalent double nécessairement\fg.  
    Dire si cette affirmation est vraie ou fausse, puis la démontrer.
    \item On souhaite obtenir un dosage plus précis. Deux élèves proposent :
    \begin{itemize}
        \item élève 1 : diluer la solution acide avant prélèvement ;
        \item élève 2 : utiliser une solution basique moins concentrée.
    \end{itemize}
    Discuter les effets de ces deux choix sur la valeur de \(V_E\) et sur la précision du titrage.
    \item \textbf{Question ouverte.}  
    Sans faire d’application numérique, proposer un protocole réaliste permettant de préparer une dilution par un facteur \(10\) d’une solution commerciale.  
    On attend une réponse rédigée, argumentée, et utilisant le vocabulaire expérimental correct.
\end{enumerate}

%===========================================================
\section{Lecture graphique d’un spectre IR et interprétation}

On étudie un liquide organique incolore inconnu \(X\).  
On a relevé expérimentalement son spectre infrarouge, représenté ci-dessous.

\begin{schemabox}
\centering
\textbf{Spectre IR simplifié du composé \(X\)}\\[0.5em]
\begin{tikzpicture}
\begin{axis}[
    width=14.5cm,
    height=6.2cm,
    xmin=4000, xmax=600,
    ymin=0, ymax=105,
    axis x line=bottom,
    axis y line=left,
    xlabel={Nombre d'onde \((\mathrm{cm}^{-1})\)},
    ylabel={Transmittance (\%)},
    xtick={4000,3500,3000,2500,2000,1700,1500,1200,1000,800,600},
    ytick={0,20,40,60,80,100},
    grid=major,
    domain=600:4000,
    samples=400
]
\addplot[thick,smooth] coordinates {
(4000,92) (3800,92) (3600,91) (3400,90)
(3250,70) (3150,45) (3050,30) (2950,28) (2850,35) (2750,48) (2650,67) (2550,84)
(2400,91) (2200,92) (2000,92) (1850,90)
(1760,82) (1710,18) (1660,82)
(1600,88) (1500,84) (1450,76) (1400,88)
(1300,90) (1250,78) (1200,64) (1150,70)
(1100,80) (1000,86) (900,90) (800,91) (700,92) (600,92)
};

% repères verticaux
\draw[dashed] (axis cs:3200,0) -- (axis cs:3200,100);
\draw[dashed] (axis cs:2500,0) -- (axis cs:2500,100);
\draw[dashed] (axis cs:1700,0) -- (axis cs:1700,100);

% petite graduation renforcée sur l'axe pour 1700
\draw[very thick] (axis cs:1700,0) -- (axis cs:1700,5);

% annotations
\node[above] at (axis cs:2850,95) {bande large};
\node[above right] at (axis cs:1700,95) {pic vers \(1700\ \mathrm{cm}^{-1}\)};

% flèche vers le minimum
\draw[->,thick] (axis cs:1820,35) -- (axis cs:1710,18);

\end{axis}
\end{tikzpicture}
\end{schemabox}

\begin{schemabox}
\centering
\textbf{Repères utiles : intervalles de bandes caractéristiques}\\[0.5em]



\medskip
\begin{tabular}{>{\bfseries}m{3.8cm} m{9.8cm}}
Alcool \(O-H\) & bande large vers \(3200\) à \(3600\ \mathrm{cm}^{-1}\) \\
Acide carboxylique \(O-H\) & très large bande vers \(2500\) à \(3200\ \mathrm{cm}^{-1}\) \\
Groupe carbonyle \(C{=}O\) & bande forte vers \(1650\) à \(1750\ \mathrm{cm}^{-1}\) \\
Liaison \(C-O\) d’un ester & bande souvent marquée vers \(1000\) à \(1300\ \mathrm{cm}^{-1}\) \\
Liaison \(C-H\) d’alcane & bande vers \(2850\) à \(3000\ \mathrm{cm}^{-1}\)
\end{tabular}
\end{schemabox}

Le composé \(X\) est supposé appartenir à l’une des familles suivantes :
\begin{itemize}
    \item alcool ;
    \item acide carboxylique ;
    \item ester ;
    \item cétone.
\end{itemize}

\begin{enumerate}[label=\textbf{\arabic*.}]
    \item À partir du \textbf{spectre IR}, relever les deux zones les plus remarquables du document expérimental.
    \item En vous aidant du \textbf{graphique des intervalles de bandes caractéristiques}, associer à chacune de ces zones une liaison ou un groupe caractéristique probable.
    \item Expliquer pourquoi la présence d’une bande forte vers \(1700\ \mathrm{cm}^{-1}\) ne suffit pas, à elle seule, à conclure que le composé est une cétone.
    \item Éliminer, en justifiant, les familles incompatibles avec le spectre.
    \item Déterminer la famille fonctionnelle la plus probable pour le composé \(X\).
    \item Expliquer pourquoi ce spectre ne permet pas forcément d’identifier \textit{de manière unique} la molécule exacte.
    \item \textbf{Question ouverte.}  
    On veut maintenant distinguer expérimentalement, au laboratoire, deux liquides incolores : un alcool et un acide carboxylique.  
    Proposer une stratégie expérimentale simple, fondée sur les connaissances du programme de Terminale, et justifier chaque étape.
\end{enumerate}

%===========================================================
\section{Diffraction : établir un modèle puis l’exploiter}

On éclaire une fente de largeur \(a\) par un faisceau laser monochromatique de longueur d’onde \(\lambda\).  
Un écran est placé à une distance \(D\) de la fente.  
On observe sur l’écran une tache centrale de largeur \(L\).

On admet que le premier minimum de diffraction est observé pour un angle \(\theta\) tel que :
\[
\theta \approx \frac{\lambda}{a}
\]
dans l’approximation des petits angles.

\begin{schemabox}
\centering
\textbf{Schéma explicatif de la diffraction par une fente}\\[0.5em]
\begin{tikzpicture}[scale=0.9]
% source
\draw[thick,->] (-4,1.5) -- (-1.2,1.5);
\node[left] at (-4,1.5) {laser};

% fente
\draw[thick] (-1,0) -- (-1,3);
\draw[line width=2.2pt,white] (-1,1.2) -- (-1,1.8);
\draw[<->] (-0.7,1.2) -- (-0.7,1.8);
\node[right] at (-0.7,1.5) {\(a\)};

% screen
\draw[thick] (5,0) -- (5,3);
\node[right] at (5,2.9) {écran};

% rays
\draw[thick] (-1,1.5) -- (5,1.5);
\draw[dashed] (-1,1.5) -- (5,2.35);
\draw[dashed] (-1,1.5) -- (5,0.65);

% angle
\draw (0.5,1.5) arc[start angle=0,end angle=8,radius=1.2];
\node at (1.7,1.8) {\(\theta\)};

% width L
\draw[<->] (5.3,0.65) -- (5.3,2.35);
\node[right] at (5.3,1.5) {\(L\)};

% D
\draw[<->] (-1,-0.5) -- (5,-0.5);
\node at (2,-0.8) {\(D\)};
\end{tikzpicture}
\end{schemabox}

\begin{enumerate}[label=\textbf{\arabic*.}]
    \item À l’aide du schéma, expliquer pourquoi on peut relier \(L\), \(D\) et \(\theta\).
    \item Montrer que, pour les petits angles :
    \[
    L \approx \frac{2\lambda D}{a}.
    \]
    \item Sans calcul numérique, indiquer comment varie la largeur \(L\) lorsque :
    \begin{itemize}
        \item la largeur \(a\) de la fente augmente ;
        \item la distance \(D\) à l’écran augmente ;
        \item la longueur d’onde \(\lambda\) augmente.
    \end{itemize}
    \item Un élève affirme : \og si on divise la largeur de la fente par \(2\), la figure de diffraction est deux fois moins large, car la lumière passe moins\fg.  
    On en pense quoi?
    \item On réalise deux expériences avec la même fente :
    \begin{itemize}
        \item expérience A : laser rouge \(a\) ;
        \item expérience B : laser vert \(2a\).
    \end{itemize}
    Comparer les largeurs des taches centrales.
    \item \textbf{}  
    On souhaite mesurer expérimentalement la largeur \(a\) d’un cheveu à l’aide d’un laser.  
    Proposer un protocole, expliquer le principe physique utilisé, préciser les mesures à effectuer, et établir la formule finale permettant d’obtenir \(a\).
\end{enumerate}

%===========================================================
\section{Projectile : modèle, trajectoire et discussion des hypothèses}

Dans tout l’exercice, on étudie le mouvement d’un projectile lancé depuis un point \(A\) de coordonnées \((0,h)\) dans un repère orthonormé \((Ox,Oz)\), avec \(Ox\) horizontal et \(Oz\) vertical vers le haut.

La vitesse initiale a pour norme \(v_0\) et fait un angle \(\alpha\) avec l’horizontale.

Dans un premier temps, on \textbf{néglige les frottements de l’air}.

\begin{schemabox}
\centering
\textbf{Schéma explicatif du lancer}\\[0.5em]
\begin{tikzpicture}[scale=0.95]
\draw[->,thick] (-0.2,0) -- (7,0) node[right] {\(x\)};
\draw[->,thick] (0,-0.2) -- (0,4.8) node[above] {\(z\)};

% trajectoire partant de A et passant au-dessus de l'obstacle
\draw[thick,domain=0:6,samples=100] plot (\x,{1.2 + 1.35*\x - 0.18*\x*\x});

\fill (0,1.2) circle (0.06);
\node[left] at (0,1.2) {\(A(0,h)\)};

\draw[->,thick] (0,1.2) -- (1.5,2.5);
\node[above right] at (1.5,2.5) {\(\vec v_0\)};

\draw (0.8,1.2) arc[start angle=0,end angle=40,radius=0.8];
\node at (1.2,1.55) {\(\alpha\)};

\draw[dashed] (0,1.2) -- (1.6,1.2);
\draw[dashed] (4.7,0) -- (4.7,2.0);
\draw[thick] (4.7,0) rectangle (4.95,2.0);
\node[right] at (4.95,1.0) {obstacle};
\end{tikzpicture}
\end{schemabox}

\begin{enumerate}[label=\textbf{\arabic*.}]
    \item Donner les coordonnées du vecteur vitesse initiale.
    \item En appliquant la deuxième loi de Newton, montrer que l’accélération du projectile est constante et déterminer ses coordonnées.
    \item En déduire les expressions de \(x(t)\) et \(z(t)\).
    \item Donner l'équation trajectoire, quel est le type de cette trajectoire ?
    \item Expliquer pourquoi le mouvement horizontal est uniforme alors que le mouvement vertical ne l’est pas.
    \item Un élève affirme : \og au point le plus haut, la vitesse du projectile est nulle puisque le projectile s’arrête avant de redescendre\fg.  
    Dire si cette affirmation est vraie ou fausse, puis justifier rigoureusement.
    \item Déterminer littéralement l’instant auquel l’altitude est maximale.
    \item En déduire littéralement la hauteur maximale atteinte.
    \item \textbf{}  
    On veut savoir si le projectile franchit un obstacle vertical situé à une distance horizontale donnée.  
    Expliquer, sans utiliser de valeurs numériques, comment décider si le projectile franchit ou non cet obstacle.  
    On attend une méthode générale, clairement rédigée.
    \item \textbf{Problème ouvert : comment décider si un modèle sans frottements est acceptable ?}  
    Dans de nombreuses situations réelles (tir d’un ballon, lancer d’une balle, jet d’eau, etc.), on utilise un modèle dans lequel on néglige les frottements de l’air.

    Rédiger une discussion argumentée répondant à la question suivante :
    \begin{center}
        \textit{Dans quelles situations ce modèle peut-il être considéré comme acceptable, et quelles sont ses limites ?}
    \end{center}

    On attend en particulier :
    \begin{itemize}
        \item une explication du sens physique de l’hypothèse \og sans frottements\fg ;
        \item des exemples de situations où cette hypothèse peut sembler raisonnable ;
        \item des exemples de situations où elle devient discutable ;
        \item une réflexion sur les grandeurs qui peuvent être modifiées par la présence des frottements (portée, durée du vol, forme de la trajectoire, etc.) ;
        \item une conclusion nuancée sur l’intérêt et les limites du modèle.
    \end{itemize}
\end{enumerate}

%===========================================================
\section{Pourquoi une sirène paraît-elle différente quand elle s’éloigne ?}

On écoute le son émis par une sirène de véhicule en mouvement.  
Un observateur immobile remarque que le son perçu lui semble plus aigu lorsque le véhicule s’approche, puis plus grave lorsqu’il s’éloigne.

On s’intéresse à cette situation d’un point de vue qualitatif, à partir des notions de signal périodique, de fréquence et de propagation d’une onde sonore.



\end{document}